% sections/resultados.tex
\section{Resultados e Discussões}

\lipsum[12-13] Os resultados obtidos demonstram a eficácia da metodologia proposta. A Tabela \ref{tab:resultados} apresenta uma comparação quantitativa entre diferentes abordagens.

\begin{table}[H]
\centering
\caption{Comparação de desempenho entre diferentes métodos}
\label{tab:resultados}
\begin{tabular}{@{}lccc@{}}
\toprule
\textbf{Método} & \textbf{RMSE} & \textbf{Tempo (s)} & \textbf{Precisão (\%)} \\ \midrule
Método Proposto & 0.0023 & 15.2 & 98.5 \\
Euler Clássico & 0.0145 & 8.7 & 92.3 \\
Runge-Kutta 4ª ordem & 0.0038 & 21.5 & 97.1 \\
Método Adaptativo & 0.0031 & 18.9 & 97.8 \\
\bottomrule
\end{tabular}
\end{table}

\subsection{Análise Quantitativa}

\lipsum[14] Os dados quantitativos revelam que o método proposto apresenta melhor desempenho em termos de precisão. A relação entre os parâmetros pode ser expressa por:

\begin{equation}
\eta = \frac{P_{out}}{P_{in}} \times 100\% = \frac{V_{out} \cdot I_{out}}{V_{in} \cdot I_{in}} \times 100\%
\label{eq:efficiency}
\end{equation}

onde $\eta$ representa a eficiência do sistema, $P_{out}$ é a potência de saída e $P_{in}$ é a potência de entrada.

\subsection{Análise Qualitativa}

\lipsum[15] Do ponto de vista qualitativo, observou-se que o comportamento do sistema é fortemente influenciado pelas condições iniciais e pelos parâmetros de controle. A matriz de transformação pode ser representada como:

\begin{equation}
\mathbf{A} = \begin{bmatrix}
a_{11} & a_{12} & a_{13} \\
a_{21} & a_{22} & a_{23} \\
a_{31} & a_{32} & a_{33}
\end{bmatrix}
\label{eq:matrix}
\end{equation}

\subsection{Discussão}

\lipsum[16-17] A análise dos resultados indica que a abordagem proposta oferece vantagens significativas em relação aos métodos tradicionais. Lorem ipsum dolor sit amet, consectetur adipiscing elit. Nullam in dui mauris. Vivamus hendrerit arcu sed erat molestie vehicula.

Proin adipiscing porta tellus, ut feugiat nibh adipiscing sit amet. In eu justo a felis faucibus ornare vel id metus. Vestibulum ante ipsum primis in faucibus orci luctus et ultrices posuere cubilia curae.
